\section{Andmed ja andmete ettevalmistus}
\label{andmed}

Uurimistöös lähtuti Eesti avaandmetest, kasutades lamapuidu kaardistamiseks Maa- ja Ruumiameti aerolaserskaneerimise (ALS) andmeid. Kõrge tuvastustäpsusega masinõppemudeli treenimiseks vajaliku andmestiku loomine vajab sihipärast uuringuala valikut, LiDAR-punktipilvedest madala taimkatte kõrgusmudelite (CHM) arvutamist ning ruumiandmete sildistamise tööriista kasutamist referentsandmete (ingl \textit{ground truth}) kogumiseks.

\subsection{Aerolaserskaneerimine (LiDAR)}

LiDAR (Light Detection and Ranging) põhineb aktiivsel kaugseirel: sensor saadab välja lühikesi laserimpulsse, mõõdab impulsi teekonna aega ja arvutab tagasipeegeldunud signaali põhjal objekti kauguse. Kui sama põhimõtet rakendatakse suurel hulgal impulssidel, tekib punktipilv, mille tihedus sõltub lennukõrgusest, lennukiirusest, impulsside sagedusest ning maapinna ja taimkatte struktuurist. Madalam lennukõrgus ja aeglasem lennu trajektoor annavad üldjuhul tihedama punktipilve, kõrgem lend aga suurema katvuse väiksema detailsusega.

Lamapuidu tuvastamise seisukohalt on oluline, et punktipilv sisaldaks piisavalt maapinnalähedasi tagasipeegeldusi, kuid samal ajal ei muutuks maapinna reljeef ja madal taimestik liiga mürarikkaks. Seetõttu ei piisa lamapuidu puhul ainult punktide olemasolust, vaid oluline on ka nende ruumiline kontsentratsioon ja fenoloogiline aastaaeg andmete kogumisel.



\subsection{Maa- ja Ruumiameti aerolaserskaneerimise kõrguspunktid}

Magistritöö aluseks on Maa- ja Ruumiameti aerolaserskaneerimise (ALS) andmestik aastatest 2017--2024. Andmeid on kogutud mitmel erineva metoodikaga, mille punktitihedus ja katvuse eesmärk erinevad. Töö kontekstis on oluline eristada riiklikke baaskaardi tootmise lende ja täpsemaid tiheasustusalade kaardistuslende, sest just nende erinevus määrab, kas andmestik sobib antud töös kasutamiseks. Tiheasustusalasid lennatake iga aastaselt kuid põhikaardi tootmiseks lennatakse umbes üks neljandik eestist korraga, nii kaetakse kogu eesti iga 4 aasta tagant. Antud töös ei sobinud kasutada suvel toimuvad metsanduslennud, sest maapinniani jõudvaid punkte vähe tänu kõrgele lennukõrgusele, lehtedele ja alustaimestikule.

Maa- ja Ruumiameti töötlusloogika sisaldab tavaliselt punktide sidumist, punktide klassifitseerimist ning kvaliteedikontrolli, et punktipilvest saaks tekitada erinevaid tooteid. Antud töös kasutati eelkõige tihedamat madallendude andmestikku.

\begin{table}[h!]
\centering
\caption{Maa- ja Ruumiameti lennukategooriate võrdlus.}
\label{tab:andmed-lennud}
\begin{tabular}{lllll}
\hline
\textbf{Lennu liik} & \textbf{Lennukõrgus} & \textbf{Punktitihedus} & \textbf{Otstarve} & \textbf{Ajaline aken} \\
\hline
Tiheasustuse / madallend & u 1200 m & 18 p/m$^2$ & tiheasustusalad & Kevad \\
Pildistamise prioriteet & u 2000 m & 3,5 p/m$^2$ &  & Kevad \\
Üle-eestiline kaardistus & u 2600 m & 2,1 p/m$^2$ & Põhikaardi tootmine & Kevad \\
Metsanduslik lend & u 3100 m & 0,8 p/m$^2$ & Metsanduslik & Suvi \\
\hline
\end{tabular}
\end{table}

\begin{figure}[h!]
\centering
\fbox{\parbox{0.92\linewidth}{\centering
\textbf{Joonis 3.1.} Maa- ja Ruumiameti andmekogum: võrdlus madallendude ja riiklike baaskaardistamise lendude vahel.\
Placeholder}}
\caption{Maa- ja Ruumiameti lennukategooriate võrdlus.}
\label{fig:mara-flights}
\end{figure}

\subsection{Uuringuala ja valimi moodustamine}

Uuringuala valik tehti kahes etapis. Esmalt kasutati Tartu linna lähedasi metsi, mis olid tänu tiheasustusala kaardistusele tihedama punktitihedusega. See oli vajalik, et arendada ja kontrollida, kas madala taimkatte kõrgusmudeli ja sildistamise töövoog töötab hästi väiksemal, kuid kvaliteetsel andmestikul. Seejärel laiendati valim üle Eesti nii, et hõlmatud oleksid erinevad metsatüübid, eri aastad ja erineva punktitihedusega kaardilehed.

Lõplik valim koosneb 29 valitud kaardilehe (1:2000) asukohast, kus valiti kõikide aastate ülelennud mis olid madala ülelennuga kogutud aastatel 2017 kuni 2024 ning mille kogupindala on 117 km$^2$. Ekspertvalim on koostatud nii, et see kataks erinevaid reljeefiga alasid, metsa- ja maastikutüüpe ning oleks kaetud tihedama ALSi andmestikuga. See tagab, et tulemused on üldistuvad kogu Eestile.

Andmestiku statistika näitab, et uuritavad Maa- ja Ruumiameti aerolaserskaneerimise andmed katavad kokku 27 km$^2$ unikaalset pindala 29 kaardilehel, kogusummas aga 117 km$^2$ (arvestades kõiki aastate ülelendude korratusi). Andmestik sisaldab kokku 4,45 miljardit laserpunkti keskmise tihedusega 37,86 pts/m$^2$ (vahemikus 11,81--81,32 pts/m$^2$). See märkimisväärselt kõrgem tihedus võrreldes standardmõõtudega peegeldab tiheasustusala lende, mida kasutati magistritöös madala taimkatte mudelite loomiseks.

\begin{figure}[h!]
\centering
\fbox{\parbox{0.92\linewidth}{\centering
\textbf{Joonis 3.2.} Eesti kaart 29 valitud kaardilehega, värvus vastavalt  punktitihedusele.\
Placeholder}}
\caption{Uuringuala valim ja kaardilehtede jaotus.}
\label{fig:study-area}
\end{figure}

\begin{table}[ht]
\centering
\begin{tabular}{ll}
\toprule
Näitaja & Väärtus \\
\midrule
LAZ-failide koguarv & 119 \\
Unikaalsed kaardilehed & 29 \\
Uurimisperiood & 2017–2024 (8 aastat) \\
Kampaania tüüp & \SI{118}{} madal, \SI{1}{} tava \\
Unikaalne pindala (kaardilehe ulatus) & \SI{27}{\kilo\meter\squared} \\
Mõõdistatud kogupindala (kõik aastad) & \SI{117}{\kilo\meter\squared} \\
Laserpunktide koguarv & \SI{4.45}{\billion} \\
Keskmine punktitihedus & \SI{37.86}{pts/\meter\squared} \\
Mediaanpunktitihedus & \SI{38.11}{pts/\meter\squared} \\
Tiheduse vahemik (min–max) & \SI{11.81}{} – \SI{81.32}{pts/\meter\squared} \\
Algandmete andmemaht & $\sim$\SI{46}{\giga\byte} \\
\bottomrule
\end{tabular}
\caption{Andmestiku ulatus ja omadused.}
\label{tbl:laz_dataset_overview}
\end{table}

\begin{table}[ht]
\centering
\begin{tabular}{lrrrrrr}
\toprule
Aasta & Faile & Kaardiruute & Pindala (\si{\kilo\meter\squared}) & Punkte (mln) & Keskmine tihedus (\si{pts/\meter\squared}) \\
\midrule
2017 & 11 & 11 & 11 & 266.7 & 24.25 \\
2018 & 17 & 17 & 17 & 553.0 & 32.53 \\
2019 & 13 & 13 & 13 & 521.5 & 40.12 \\
2020 & 18 & 18 & 18 & 723.2 & 40.18 \\
2021 & 14 & 14 & 14 & 527.6 & 37.68 \\
2022 & 26 & 26 & 24 & 1040.5 & 42.05 \\
2023 & 12 & 12 & 12 & 481.4 & 40.12 \\
2024 & 8 & 8 & 8 & 338.9 & 42.36 \\
\midrule
Total & 119 & — & 117 & 4452.7 & 37.86 \\
\bottomrule
\end{tabular}
\caption{Per-year breakdown of coverage, files, and density.}
\label{tbl:laz_per_year}
\end{table}

Kaardiruutude detailne koondtabel asub lisas \ref{lisa:kaardiruudud}.

\subsection{CHM ja eeltöötlus}

Lamapuidu märgistamise jaoks tekitati madal taimkattemudel ( Canopy Height Model -- CHM) järgmise töövoo järgi: punktipilvest maapinnamudel, arvutati kõrgus maapinnast (Height Above Ground -- HAG), filtreeriti välja punktid kõrgusvahemikust 0--1,3 m ning rasteriseeriti maksimaalne kõrgus 20 cm piksliga rastriks. 1,3 m piir on siin kasutatud ka varasemates töödes \cite{jarron_detection_2021, joyce_detection_2019}: see jätab välja kõrgema ja tõenäolisemalt elusa taimestiku ning vähendab müra, kuid hoiab alles lamapuiduga seotud maapinnalähedastest struktuuridest ja võimaldab kasutada agrigeerimiseks maksimaalset kõrgust taimkatte mudeli jaoks ilma olulist infot varjamata. Selline lähenemine hoiab CHM-i semantika lihtsana ja vähendab kõrgema taimestiku mõju. Mõju ei kao kuna igihaljad puud vähendavad laserimpulsside läbilastvust ja vähendab maapinna lähedastese kididesse jõudmist seega tekitavad varje ja muudavad punktitiheduse väikesel alal väga ebaühtlaseks. 

Andmete rasteriseerimisel rakendati piksli suurust 20 cm, omistades igale piksli ruudule selles sisalduva punktipilve maksimaalse kõrgusväärtuse. Resolutsiooni valik on metoodiliselt kaalutletud otsus, mis peab tasakaalustama punktipilve tihedust, andmestiku detailset katvust ning võimekust eristada kõrvuti asetsevaid üksikobjekte. 

Antud uuringus valitud 20 cm piksli suurus on otseses seoses tuvastatavate objektide füüsiliste mõõtmetega. Statistilise metsainventuuri (SMI) 2024. aasta välitööde juhendi kohaselt loetakse lamapuidu mõõtmisel künniseks tüve läbimõõt alates 15 cm \cite{SMI2024}. Samas rõhutab erialane kirjandus, et elurikkuse seisukohalt on kriitilise tähtsusega just jämedam lamapuit — paljud nõudlikud ja ohustatud saproksüülsed liigid eelistavad või vajavad elupaigana tüvesid, mille läbimõõt on vähemalt 20 cm \cite{Vilhelmsson2013}. Seega on 20 cm piksli suurus optimaalne lävi, mis vastab minimaalsele ökoloogiliselt olulisele objekti suurusele, mida käesoleva analüüsi raames soovitakse tuvastada, tagades samas piisava ruumilise täpsuse objektide kuju säilitamiseks.

\begin{figure}[h!]
\centering
\fbox{\parbox{0.92\linewidth}{\centering
\textbf{Joonis 3.3.} Okaspuude varjud maapinnale jõudvate punktide peal, mille tulemusena väheneb punktitihedus tihedates metsades.\
Placeholder}}
\caption{Okaspuude varjude mõju punktipilve tihedusele metsaalades.}
\label{fig:shadow-density-effect}
\end{figure}

\begin{figure}[h!]
\centering
\fbox{\parbox{0.92\linewidth}{\centering
\textbf{Joonis 3.4.} Toor-punktipilv -> DTM-i eemaldamine -> normaliseeritud kõrgused -> filtreeritud punktid (0--1,3 m) -> rasterpilt (20 cm).\
Placeholder}}
\caption{CHM-i loomise ja eeltöötluse voogdiagramm.}
\label{fig:chm-flow}
\end{figure}

\begin{figure}[h!]
\centering
\fbox{\parbox{0.92\linewidth}{\centering
\textbf{Joonis 3.5.} Näide ühest objektist: punktipilv (3D) versus genereeritud 20 cm CHM raster.\
Placeholder}}
\caption{Lamapuu objekti esitus punktipilves ja CHM-is.}
\label{fig:chm-example}
\end{figure}

Selle töötluse tulemuseks saadud detailsed ja madalad taimkattemudelid on alus nii sidistamisele kui ka mudeli sisendile. Eesmärk ei ole üksnes objektide nähtavaks tegemine, vaid ka erinevate aastate ja alade vahel võrdleva, ühtse medoodika alusel kujutise loomine.

Sildistamise ja hilisema mudeli sisendi jaoks jagati suuremõõtmelised ($1 \times 1$ km) TIF-failid väiksemateks, $128 \times 128$ piksli suurusteks pildikildudeks (ingl \textit{tiles}). Kuna rasteri resolutsiooniks oli eelnevalt määratud 20 cm piksli kohta, vastab ühe tiili maastikulisele ulatusele $25,6 \times 25,6$ meetrit. Selline suuruse valik ei olnud juhuslik, vaid tugines mitmele disainikaalutlusele:
\begin{itemize}
    \item \textbf{Töövoo kiirus:} $128 \times 128$ piksline pilt laeb kiiresti ja kuvatakse brauseris sujuvalt, muutes suuremahulise sildistamisprotsessi inimoperaatorile visuaalselt hoomatavaks ning mugavaks.
    \item \textbf{Arhitektuuriline paindlikkus:} Küljepikkus 128 kujutab endast kahe astet ($2^7$), mis on ideaalne ja valdavalt standardnõue konvolutsiooniliste närvivõrkude ja süvaõppearhitektuuride mastaapimisel. Kuna andmete ettevalmistamisel polnud lõplik mudel veel fikseeritud, tagas selline mõõt universaalse sobivuse erinevate arhitektuuridega.
    \item \textbf{Kontekstuaalne ülevaade:} $25,6$ meetri pikkune ala on piisavalt lai, et tavapärasesse metsa maha langenud eesti lamapuidu objekt, mis jääb üldiselt samasse suurusjärku, oleks pildil nähtav tervikuna koos seda ümbritseva reljeefiga.
    \item \textbf{Ülekate servaalade kompenseerimiseks:} Piltide lõikamisel rakendati igas suunas 50\% ülekatet (samm 64 pikslit). See tagab, et objekt, mis satuks muidu pildi servale või lõigataks pooleks, jääb vähemalt ühel tiilil pildi keskossa. Pidev ja korduv keskse asetusega treeninginfo võimaldab mudelil objekte usaldusväärsemalt tuvastada.
\end{itemize}
